% =====================================================================
% CUMCM 论文主骨架（comp_cumcm 模板 · 统一工作流 S6/S8 使用）
% 源模板：templates/comp_cumcm/latex/（cumcmthesis.cls 为开源模板 v2.6，
%         出处：Modex-MH-Agent 安装包内置竞赛模板，社区许可使用）
% 编译：  python scripts/build_latex.py paper/latex/main.tex   （xelatex 两遍）
% 装配：  把 paper/sections/*.md 转为同名 *.tex 放入本目录后填写 TODO 项。
% 版式合同（CLAUDE.md）：西文与正文数字 Times New Roman；一级标题 16pt（三号）；
%         目录条目比正文小一档；每张图图下自明标题 + 紧随其后的读图段落。
% =====================================================================
\documentclass[withoutpreface]{cumcmthesis}
% withoutpreface 去掉承诺书/编号页，提交时按当年通知换回（或使用官方封面页）。

% ---- 字体：西文与数字 Times New Roman ----
\setmainfont{Times New Roman}
\setsansfont{Times New Roman}
\setmonofont{Consolas}

% ---- 一级标题 16pt（三号）----
\usepackage{titlesec}
\titleformat{\section}{\zihao{3}\bfseries\centering}{第\,\thesection\,章}{0.5em}{}
\titlespacing{\section}{0pt}{1.2em}{0.8em}

% ---- 图表与引用 ----
\usepackage{booktabs}
\usepackage{multirow}
\usepackage{caption}
\captionsetup{font=small,labelfont=bf}

% =====================================================================
% 摘要：背景/总体框架一段 + 每问独立一段 + 综合边界一段 + 关键词
% =====================================================================
\begin{document}

\title{题名待填：<题目短名> 的建模与求解}
\author{队号：\mcm@tokens@baominghao}  % TODO 提交时填写
\date{\today}

\maketitle

\begin{abstract}
% TODO 段落1：背景与总体框架（数据、方法体系、验证链路、递进框架）
% TODO 段落2-5：问题一至问题四各独立成段（方法与带单位核心输出）
% TODO 段落6：综合边界与适用范围
\keywords{关键词一；关键词二；关键词三；关键词四；关键词五}
\end{abstract}

\tableofcontents

% =====================================================================
% 第一章 问题背景与重述
% =====================================================================
\section{问题背景与重述}
\input{chapters/01-problem-analysis.tex}   % TODO 由 paper/sections/01-problem-analysis.md 转换

% =====================================================================
% 第二章 问题分析
% =====================================================================
\section{问题分析}
\input{chapters/02-analysis.tex}           % TODO 转换/撰写

% =====================================================================
% 第三章 模型假设
% =====================================================================
\section{模型假设}
\input{chapters/03-assumptions.tex}        % TODO 转换/撰写

% =====================================================================
% 第四章 符号说明
% =====================================================================
\section{符号说明}
\input{chapters/04-symbols.tex}            % TODO 转换/撰写（符号表）

% =====================================================================
% 第五章 模型的建立与求解（四问为二级标题，三级标题点明本问对象）
% =====================================================================
\section{模型的建立与求解}

\subsection{问题一：<本问对象>}
% 三级标题按职责锚点 + 本问对象，例如：
% \subsubsection{任务转化与机理假设}
% \subsubsection{模型建立与推导（模型名、变量）}
% \subsubsection{求解流程与赛马比较（MLR vs Hill）}
% \subsubsection{核心结果与读图解释}
% \subsubsection{方案选择与模型诊断}
\input{chapters/q1.tex}                    % TODO 由 q1-{modeling-process,results,model-argumentation}.md 合并转换

\subsection{问题二：<本问对象>}
\input{chapters/q2.tex}

\subsection{问题三：<本问对象>}
\input{chapters/q3.tex}

\subsection{问题四：<本问对象>}
\input{chapters/q4.tex}

% =====================================================================
% 第六章 模型评价与应用边界
% =====================================================================
\section{模型评价与应用边界}
\input{chapters/07-evaluation.tex}         % TODO 转换/撰写

% =====================================================================
% 第七章 结论
% =====================================================================
\section{结论}
\input{chapters/08-conclusion.tex}         % TODO 转换/撰写

% =====================================================================
% 参考文献
% =====================================================================
\bibliographystyle{unsrt}
\nocite{*}
\bibliography{refs}                        % TODO refs.bib 或直接 thebibliography

% =====================================================================
% 附录（可选；正文与附录分别计数）
% =====================================================================
%\appendix
%\section{主要代码与补充图表}

\end{document}
